\chapter{Fundamentos teóricos}\label{ch:theoretical_background}

% ============================================================
%  BLOQUE A — CONTEXTO DINÁMICO
% ============================================================

\section{Ciclones extratropicales en el Atlántico Sur}\label{sec:tb_fundamentos}

La teoría clásica, desarrollada a partir de los primeros estudios sinópticos de ciclones de latitudes medias, los definió como vórtices de núcleo frío que evolucionan a partir de una onda sobre un frente de discontinuidad térmica entre masas de aire frío y cálido \citep{bjerknes1922life}.
Este modelo fue revisado décadas después por \citet{shapiro1990life}, cuyo marco estructural se desarrolla en la Sección~\ref{subsec:tb_modelos}.
En el Hemisferio Sur, \citet{hoskins2005new} generaron la climatología y confirmaron la concentración de la ciclogénesis en las zonas baroclínicas.
En el Atlántico Sur, \citet{sinclair1995climatology} resaltó que la ciclogénesis está influenciada por los gradientes térmicos oceánicos y la orografía, mientras que \citet{inatsu2004zonal} demostraron que la asimetría zonal del \textit{storm track} está organizada por ondas estacionarias troposféricas, moduladas regionalmente por la topografía de los Andes, cuya interacción con el flujo de niveles bajos también favorece la propagación de ondas baroclínicas sobre Sudamérica \citep{vera2002cold}.
Si bien la inestabilidad baroclínica provee el entorno favorable para la génesis, procesos diabáticos asociados a la liberación de calor latente modulan la evolución e intensidad del sistema, como concluyó \citet{coutodesouza2024thesis}. 
Esto concuerda con experimentos de sensibilidad numérica de \citet{reboita2022from}, donde la supresión de los flujos de calor sensible y latente alteró la estructura térmica y la intensidad del vórtice, efecto también reportado por \citet{gozzo2013air}.

% ============================================================
%  BLOQUE B — DETECCIÓN Y SEGUIMIENTO
% ============================================================
\section{Detección de ETC}\label{sec:tb_vorticidad}

\subsection{Marco de referencia semi-lagrangiano}\label{subsec:tb_lagrangiano}
La variabilidad estructural de los ETCs del Atlántico Sur requiere un marco de referencia común que no diluya la señal en promedios geográficos ni incorpore ruido de sistemas adyacentes.
El enfoque euleriano superpone campos de variables, dificultando observar asimetrías estructurales durante la evolución. El enfoque lagrangiano evita esta dilución centrando cada ETC en su núcleo de vorticidad con coordenadas relativas $(\Delta x, \Delta y)$, preservando la organización espacial y aislando su señal dentro de un dominio que se desplaza con la trayectoria.
Un dominio de $20^{\circ} \times 20^{\circ}$ se empleó en el Atlántico Sudoccidental por \citet{cardoso2022synoptic}. Este enfoque lagrangiano ha demostrado utilidad a escala global: \citet{laurila2021characteristics} muestran que permite construir composites en el Hemisferio Norte, mientras \citet{priestley2022improved} lo aplicaron en ambos hemisferios para cuantificar la estructura e intensidad de los vientos, \citet{pepler2020dimensional} lo emplearon con un radio de $10^{\circ}$ en el sureste de Australia para estudiar la severidad de vientos y lluvia, y \citet{dacre2012atlas} lo utilizaron con un radio de $20^{\circ}$ en el Atlántico Norte para construir un atlas de compuestos que documenta sistemáticamente la estructura y evolución del ciclón.

\subsection{Vorticidad relativa en 850 hPa}\label{subsec:tb_zeta}
En el enfoque de seguimiento lagrangiano, la vorticidad relativa en 850 hPa ($\zeta_{850}$) se ha establecido como la variable para identificar y rastrear ciclones extratropicales.
En el Hemisferio Sur estos sistemas se desarrollan inmersos en un flujo zonal de los oestes que los traslada con rapidez; por ello, en sus etapas iniciales la caída de presión suele quedar enmascarada, presentándose como vaguadas abiertas en lugar de centros cerrados \citep{sinclair1994objective}, lo que vuelve a la presión al nivel del mar (MSLP) una métrica inadecuada en la detección temprana.
Para resolver este problema, \citet{sinclair1997objective} recomendaron utilizar la vorticidad relativa en niveles bajos en lugar de la presión al nivel del mar; en esta tesis se adopta específicamente $\zeta_{850}$.
Esta variable se ha usado objetivamente en climatologías del Hemisferio Sur completo \citep{padilhareinke2024objective} y en estudios regionales del Atlántico Sur \citep{gramcianinov2019properties}, además en el HN por \citet{corner2025classification}, y segmentación objetiva del ciclo de vida por \citet{coutodesouza2024new}, que emplea $\zeta_{850}$ como variable base para delimitar las cuatro fases mediante detección de picos y valles en su serie temporal.

\begin{equation}\label{eq:vorticidad_relativa}
\zeta = \mathbf{k} \cdot (\nabla \times \mathbf{V}) = \frac{\partial v}{\partial x} - \frac{\partial u}{\partial y},
\end{equation}

\noindent donde $u$ y $v$ son las componentes zonal y meridional del viento horizontal, respectivamente.
A 850 hPa, $\zeta$ filtra el flujo zonal de traslación y captura la rotación ciclónica intrínseca inducida por forzantes de mesoescala \citep{hoskins2005new}, señal que precede a la manifestación en el campo de presión durante las etapas iniciales \citep{sinclair1994objective}, criterio adoptado como base de detección en los principales algoritmos de rastreo y bases de datos objetivas.
  
\subsection{Estratificación}\label{subsec:tb_poblaciones}
La intensidad dinámica se cuantifica mediante el mínimo de vorticidad relativa alcanzado durante la trayectoria completa, $\zeta_{850}^{\min}$. 
Esta estratificación se justifica mediante \citet{priestley2022improved}, quienes seleccionaron el 10\% de ciclones con mayor vorticidad ciclónica en su pico de intensidad.
La pertinencia de esta estratificación, independientemente de la métrica de intensidad empleada, encuentra respaldo adicional en \citet{andrade2024composite}, quienes muestran que categorías con mayor tasa de profundización exhiben comportamiento estructural coherente y diferenciado, con mayores flujos de calor superficiales. 
Esto sugiere que los sistemas más intensos, ya sea en términos de vorticidad o de presión, presentan una organización física distintiva que justifica tratarlos como subconjunto separado.
Se adopta un criterio percentílico ($p90$) para aislar los sistemas más intensos. 
En el Hemisferio Sur, donde la vorticidad ciclónica es negativa, seleccionar p90 en valor absoluto equivale al $\zeta_{850}^{\min}$ más negativo, osea el más intenso.
Para garantizar comparabilidad estructural, la homogeneidad evolutiva requiere ciclo de vida completo como la secuencia de cuatro fases evolutivas (Sección~\ref{subsec:tb_fases}) sin re-intensificaciones que introduzcan ruido. 
Este criterio retiene entre el 43\% y el 52\% de los sistemas por región (Tabla~\ref{tab:conteo_ciclones}), a partir de una versión actualizada de la base de trayectorias desarrollada originalmente por \citet{coutodesouza2024new}.
Los filtrados específicos se detallan en la Sección~\ref{sec:poblaciones_estudio}.
De lo anterior emergen dos estratos de análisis complementarios:
\begin{itemize}
\item Población Global (\textit{Full-Life-Cycle Set}): conjunto de referencia integrado por ciclones de ciclo de vida completo que representan el comportamiento medio de la región.
\item Subconjunto p90 (\textit{High-Intensity Set}): estrato de alta intensidad dinámica extraído del percentil superior del valor absoluto de $\zeta_{850}^{\min}$ dentro de la Población Global.
\end{itemize}
Esta estratificación dual responde a si una mayor intensidad del vórtice se asocia con cambios sistemáticos en la organización espacial en superficie. 
Los sistemas $p90$ mantienen una evolución temporal coherente con la Población Global (PG), alcanzando su máxima intensidad predominantemente durante la fase de madurez (Tabla~\ref{tab:fase_max_intensidad}).

% ============================================================
%  BLOQUE C — REGIONES DE CICLOGÉNESIS
% ============================================================

\section{Regiones de ciclogénesis}\label{sec:tb_regiones}
En el Atlántico Sudoccidental, \citet{gramcianinov2019properties} encontraron que la interacción entre el flujo de los oestes, la topografía de los Andes \citep{gan1994influence} y la Temperatura Superficial del Mar define tres regiones de génesis con características físicas diferenciadas, delimitación que \citet{reboita2010south} habían establecido previamente empleando también la vorticidad relativa (aunque en otro nivel atmosférico) para identificar estos focos en Sudamérica.

\subsection{Sur-sureste de Brasil (SBR)}\label{subsec:tb_sbr}
La región SBR (referida como RG1 en \citet{reboita2022from}) constituye uno de los centros de acción preferenciales para la ciclogénesis extratropical en el litoral brasileño.
Los ciclones pueden iniciarse por inestabilidad baroclínica asociada a la advección de aire cálido en niveles bajos y una vaguada en niveles medios-altos \citep{gozzo2014subtropical}, pero su intensificación está fuertemente modulada por los flujos turbulentos de calor sensible y latente en la capa límite \citep{gozzo2013air}, que organizan la convergencia de humedad e intensifican el sistema en su sector cálido. 
El análisis energético de \citet{coutodesouza2024thesis} mostró que la contribución diabática a la energía del sistema en SBR se eleva durante el invierno, a diferencia de ARG, donde esta contribución resulta la menor de las tres regiones en ambas estaciones.
\citet{gramcianinov2024early} demostraron además que, durante la ciclogénesis temprana en SBR, los modelos de alta resolución logran representar las estructuras de mesoescala de advección cálida y convergencia de humedad en capas bajas.

\subsection{Cuenca del Río de la Plata (LPB)}\label{subsec:tb_lpb}
En la región de LPB, \citet{gramcianinov2024early} resaltaron que la ciclogénesis se ve favorecida por la interacción entre la Alta Subtropical del Atlántico Sur y el Chorro de Capas Bajas de Sudamérica, que intensifica el transporte de humedad hacia el sistema, y confirmaron mediante modelos regionales que la convergencia de humedad en capas bajas constituye el mecanismo dominante de intensificación.
\citet{coutodesouza2024thesis} encontró que este desarrollo depende del transporte de humedad del Chorro de Capas Bajas de Sudamérica sobre el flanco oriental de los Andes y del Anticiclón Subtropical del Atlántico Sur hacia la costa sureste, lo que se traduce en la mayor contribución diabática a la energía del sistema entre las tres regiones, de forma consistente en ambas estaciones.
En línea con este forzante de humedad, \citet{reboita2010south} mostraron que la región se caracteriza por una alta frecuencia de sistemas ciclónicos y por su interacción con el transporte de humedad desde latitudes tropicales, patrón que \citet{laureanti2024relation} documentan también, mediante compuestos propios, para el centro-este de Brasil, donde el desplazamiento hacia el norte del SALLJ coincide con una circulación ciclónica anómala de niveles bajos que intensifica el transporte de humedad amazónica.

\subsection{Argentina (ARG)}\label{subsec:tb_arg}
El sector austral ($>40^\circ$S) responde a un régimen dinámico clásico dominado por la baroclinicidad de gran escala, a diferencia del forzamiento termodinámico de superficie que caracteriza a SBR y LPB.
\citet{hoskins2005new} mostraron que la ciclogénesis en esta banda se concentra al este de los Andes, patrón que \citet{coutodesouza2024thesis} corroboró empíricamente en su propia climatología de trayectorias para la región.
\citet{sinclair1995climatology} muestra que la ciclogénesis frente a la costa de Argentina ocurre todo el año, a diferencia de otras zonas de génesis a sotavento cuya actividad depende de gradientes térmicos oceánicos estacionales.
\citet{mendes2010climatology} identificaron en esta región un clúster de trayectorias (\textit{storm-track}) propio dentro del sector sudamericano.
\citet{gramcianinov2024early} mostraron además que los controladores físicos de ARG presentan tendencias proyectadas comparativamente débiles hacia el futuro, a diferencia de la intensificación de humedad proyectada para esas regiones.
% ============================================================
%  BLOQUE D — ESTRUCTURA INTERNA
% ============================================================

\section{Estructura interna}\label{sec:tb_evolucion}

\subsection{Modelo conceptual de estructura interna}\label{subsec:tb_modelos}

Los ciclones extratropicales intensos del Atlántico Sur suelen seguir el modelo Shapiro-Keyser propuesto por \citet{shapiro1990life}, distinguido por cuatro rasgos diagnósticos: la fractura del frente frío (\textit{frontal fracture}), el frente curvado (\textit{bent-back front}), la estructura en T (\textit{T-bone}) y el aislamiento de una masa de aire cálido en el núcleo (\textit{warm seclusion}). 
Este modelo captura la forma que adopta el sistema, pero es el modelo de cintas transportadoras el que explica cómo esa forma se traduce en la distribución de viento y precipitación en superficie. 
El WCB, definido originalmente como la corriente ascendente de aire cálido y húmedo que se origina en el sector cálido del ciclón \citep{carlson1980airflow}, y su respectivo chorro de bajo nivel se organizan en torno al vórtice \citep{browning1986conceptual}, y el ascenso y giro anticiclónico de este flujo cálido modula las estructuras de mesoescala, concentrando las bandas de precipitación y el cizallamiento del viento por delante del frente frío \citep{coll2024lagrangian}.
El WCB se origina en el lado \textit{equatorward} del centro y asciende sobre el frente cálido, mientras que el CCB (frío y denso) ocupa el lado  \textit{poleward}, envolviendo ciclónicamente la cabeza nubosa \citep{dacre2023climatology}. 

Esta organización frontal del WCB es la que vincula su posición con los máximos de precipitación \citep{catto2015fronts}.
En el HN, el chorro del WCB domina el cuadrante sureste del ciclón y concentra allí los vientos más intensos del compuesto \citep{eisenstein2023identification, catto2010climate}, mientras que este flujo del sector cálido resulta sistemáticamente más débil en el Hemisferio Sur que en el Hemisferio Norte \citep{hodges2011comparison}.
El acoplamiento entre WCB y CCB explica esta asimetría de viento y precipitación: el CCB gana vorticidad potencial al desplazarse por el lado frío del frente curvado, bajo la región de máxima liberación de calor latente del WCB ascendente, lo que localiza el chorro de bajo nivel más intenso en el extremo del frente curvado \citep{schemm2014linkage}, patrón confirmado en ambos hemisferios \citep{priestley2022improved} y documentado en un caso real de Japón, donde la banda de precipitación estratiforme se extendió hacia el este del centro sobre el CCB \citep{sawada2021heavy}.

Bajo una simetría especular respecto al ecuador, esta organización debería invertirse únicamente en el eje norte-sur al pasar al Hemisferio Sur, donde el sentido de giro ciclónico es horario en lugar de antihorario, mientras que la disposición este-oeste debería mantenerse, al no depender de la rotación planetaria sino de la orientación del sistema respecto a su desplazamiento. Esta expectativa aún no cuenta con un estudio de viento vectorial que la confirme para el Atlántico Sur, por lo que se adopta aquí como hipótesis orientadora y no como resultado establecido.

Una intrusión seca (DI) descendiendo sobre el WCB puede generar un tercer máximo de precipitación, de carácter convectivo y localizado cerca del centro del ciclón, distinto de la banda estratiforme oriental ya descrita \citep{sawada2021heavy}.
Durante la oclusión, una corriente conocida como \textit{trowal} (\textit{trough of warm air aloft}), que asciende ciclónicamente dentro del cuadrante ocluido, reorganiza la precipitación en una banda periférica arqueada que ya no coincide con el núcleo de viento, patrón documentado en el HN \citep{martin1999forcing, naud2025lifecycle}.
El \textit{Sting Jet} (SJ), estructura de mesoescala del modelo Shapiro-Keyser, desciende desde niveles medios dentro de esa misma intrusión seca, por encima del CCB, produciendo los vientos más intensos en la zona de fractura frontal, distinta del WCB y del CCB, según casos documentados en el HN \citep{clark2005sting, clark2018sting, martinez2014cold}.

\subsection{Asimetría espacial de las variables}\label{subsec:tb_kde_asimetria}

Más allá de la asimetría espacial en superficie que aquí se analiza, cerca del 70\% de los ETCs en el Atlántico Sur presenta también una estructura térmica asimétrica en la vertical \citep{conrado2024cyclone}, un tipo de asimetría distinto pero que sugiere que la heterogeneidad estructural de estos sistemas opera en múltiples dimensiones.
\citet{mcerlich2023extremes} mostraron que la coma de precipitación rota en sentido horario, ciclónico en el HS, confirmando la inversión del sentido de circulación en relación a ciclones del HN.
Esta inversión, sin embargo, aplica al sentido de giro y geometría del WCB, no a la magnitud.
\citet{kodama2019perspective} mostraron que los ETCs oceánicos del HS precipitan sistemáticamente menos que sus contrapartes del HN (7.7 frente a 10.7~mm/día en promedio espacial), diferencia atribuible a la menor disponibilidad de humedad y no a una dinámica distinta.
Estudios del HN documentan que el WCB asciende típicamente hacia el polo con un giro hacia el este \citep{blanchard2021warm}; por efecto espejo, en el Hemisferio Sur concentraría la humedad en el sector sureste del sistema.
\citet{oertel2021warm} demostraron que el WCB frecuentemente alberga núcleos de convección intensa embebida que generan tasas de precipitación extrema, asimetría corroborada empíricamente por \citet{naud2020evaluation} y \citet{field2007precipitation}, quienes mostraron, este último incluyendo al Atlántico Sur en su análisis, que la precipitación en los ciclones oceánicos se concentra en el sector oriental/sureste en forma de coma, mientras el flanco occidental permanece subsidente y seco.
En el Atlántico Sur, \citet{gozzo2013air} documentaron que la disponibilidad de humedad superficial condiciona directamente la ubicación de la precipitación. 
Con flujos de calor activos, esta se distribuye ampliamente en torno al centro del ciclón y el frente cálido, mientras que su ausencia la debilita y concentra en el sector sureste/este del sistema.

El campo de viento en superficie presenta una asimetría que evoluciona a lo largo del ciclo de vida.
\citet{rudeva2011composite} mostraron, mediante compuestos centrados en el ciclón, que el máximo de viento se ubica al sur del centro y se intensifica marcadamente durante la etapa de profundización (de 8--9 a 14--16~m/s).
El acoplamiento WCB--CCB descrito arriba concentra los vientos más intensos en el extremo del frente curvado.
\citet{catto2016extratropical} recogen en su síntesis esta heterogeneidad en tres núcleos espacialmente distintos, WCB, CCB y sting jet, denominados WJ, CJ y SJ, respectivamente, cada uno dominante en una fase distinta del ciclo de vida, definida allí por la tasa de profundización y la estructura frontal del sistema, y no por la señal de vorticidad empleada en esta tesis.
En el Atlántico Sur, \citet{cardoso2022synoptic} documentó que los vientos máximos tienden a concentrarse en sectores específicos, como el este y sureste relativo al centro del sistema, y que las ráfagas más intensas suelen vincularse a características de mesoescala cuya posición relativa al centro de vorticidad varía a medida que el ciclón madura \citep{eisenstein2023identification}.
\citet{gentile2023observed} mostraron que, integrado sobre el ciclo de vida, el viento dañino del sector frío (asociado al CCB) es 3--4 veces más frecuente que el del sector cálido (WCB).
Esta heterogeneidad explica por qué la intensidad de los máximos de viento no puede capturarse mediante una única métrica central \citep{corner2025classification}.
Este mismo patrón se repite en la costa SE de Brasil, donde \citet{cardoso2020wind} lo asociaron a la mayor densidad de ciclones subtropicales de la región, \citet{dejesus2021multimodel} confirmaron que el sector cálido (NE en el HS) concentra los vientos en la etapa temprana de desarrollo del ciclón, y \citet{gramcianinov2024early} localizaron los vientos del noroeste más intensos al este del centro en SBR y LPB, reforzados en el cuadrante noreste de LPB por su combinación con el Chorro de Capas Bajas de Sudamérica.

\subsection{Fases del ciclo de vida}\label{subsec:tb_fases}
Para caracterizar la evolución temporal del viento y la precipitación, este estudio adopta la segmentación del ciclo de vida en cuatro fases discretas mediante la metodología objetiva del paquete computacional \textit{CycloPhaser} \citep{deSouza2025CycloPhaser}, aplicado por \citet{coutodesouza2024new}. 
Este algoritmo utiliza la tasa de cambio de la vorticidad relativa $\zeta_{850}$, conectando así directamente con el fundamento de detección establecido en la Sección~\ref{subsec:tb_zeta}, para definir los siguientes regímenes estructurales. 
Cabe notar que cada fase no es homogénea internamente. Los instantes próximos a sus transiciones comparten rasgos de la etapa adyacente, por lo que la condición estructural más representativa de cada régimen corresponde al núcleo temporal de la fase, alejado de ambos bordes de transición.

\begin{enumerate}
    \item \textit{Incipiente:} Corresponde a la aparición inicial del núcleo de vorticidad organizada.
    En el contexto regional, esta fase suele asociarse a la interacción entre una vaguada en altura y forzantes de superficie, como la orografía de los Andes o gradientes térmicos costeros.

    \item \textit{Intensificación:} Período caracterizado por el rápido crecimiento de la vorticidad ciclónica.
    Esta etapa está impulsada por una retroalimentación entre la inestabilidad dinámica y los flujos de calor latente y sensible \citep{gozzo2013air, reboita2022from}.

    \item \textit{Madurez:} Etapa donde el sistema alcanza su intensidad máxima (pico de vorticidad) y su mayor extensión horizontal. El radio de los ciclones en superficie tiende a incrementarse sistemáticamente al evolucionar hacia la madurez en ambos hemisferios, como demostró \citet{simmonds2000size}, lo que exige que el dominio de análisis lagrangiano sea suficientemente amplio para capturar la circulación completa del sistema.
    Es precisamente durante esta fase donde los ciclones oceánicos intensos consolidan la evolución estructural documentada por \citet{shapiro1999bridge}. La deformación del flujo genera la fractura del frente frío y aísla una masa de aire cálido en el núcleo del vórtice (\textit{warm seclusion}).
    Al completarse este proceso, las huellas de precipitación y viento quedan espacialmente desacopladas.

    \item \textit{Decaimiento:} Fase final marcada por la disminución progresiva de la magnitud de $\zeta_{850}$, reflejo del debilitamiento dinámico del vórtice \citep{coutodesouza2024thesis}.
\end{enumerate}

% ============================================================
%  BLOQUE E — ACOPLAMIENTO ESTRUCTURAL Y HERRAMIENTAS
% ============================================================

\section{Acoplamiento estructural}\label{sec:tb_compound_def}

Los campos de viento y precipitación en superficie no son independientes, están físicamente acoplados a través de la estructura tridimensional del ciclón, de modo que sus máximos pueden ocurrir en sectores próximos o solapados durante determinadas fases del ciclo de vida.
Viento y precipitación son, además, las dos variables de superficie más empleadas para caracterizar los impactos de un ciclón, por lo que su análisis conjunto es el punto de partida natural para caracterizar la morfología del sistema, en línea con otros enfoques que analizan la co-ocurrencia de extremos de superficie \citep{zscheischler2018future}. 
Esta co-ocurrencia, que \citet{chen2025characteristics} cuantifican en la mayoría de los extremos de viento y precipitación asociados a ETC, es consistente con el acoplamiento dinámico entre el WCB y el CCB descrito en la Sección~\ref{subsec:tb_modelos}, lo que justifica el análisis integrado de ambas variables mediante MEOF para caracterizar la morfología ciclónica.
No obstante, ERA5 presenta sesgos documentados en la representación de los extremos más intensos \citep{chen2024evaluation}, cuya magnitud y consecuencias para el subconjunto $p90$ se discuten en la Sección~\ref{sec:datos_era5}.

\section{Herramientas de caracterización}\label{sec:tb_herramientas}
Caracterizar el viento y precipitación de los ETCs requiere integrar múltiples eventos en una representación coherente. 
Las técnicas elegidas responden a propiedades intrínsecas de estos campos: asimetría, no normalidad y covariabilidad física.

\subsection{ETCs Compositing}\label{subsec:tb_compositing}
La técnica de compositing centrado en el sistema consiste en superponer los campos de distintos eventos en coordenadas relativas al centro de cada ETC, preservando las asimetrías estructurales, bajo el supuesto de que los ciclones de una misma región y fase comparten un patrón estructural estadísticamente representativo que el promedio de sus campos lagrangianos revela.
Esta lógica de alineación, que extrae los campos dentro de un radio fijo respecto al centro, ha sido empleada en la evaluación de la estructura ciclónica en modelos climáticos, como lo hizo \citet{catto2010climate} al comparar la representación de ETCs entre un modelo climático de alta resolución y el reanálisis ERA-40.
Su validez metodológica fue reafirmada por \citet{priestley2022improved}, quienes aplicaron compositing para aislar estructuras representativas de las corrientes en aire cálida (WCB) y fría (CCB), cuantificando así sesgos sistemáticos de intensidad del viento entre generaciones de modelos climáticos frente a ERA5.
Más allá de la estructura promedio estática, la técnica permite rastrear cómo los campos se configuran, desplazan y contraen a lo largo del ciclo de vida del sistema, al particionar los composites por fase de desarrollo.
\citet{mcerlich2023extremes}, aplicando esta estrategia a ETCs del Hemisferio Sur mediante ERA5, mostraron que la extensión espacial de la precipitación extrema se restringe hacia el centro del ciclón a medida que aumenta el umbral de intensidad.
En el Atlántico Sur en particular, \citet{gramcianinov2019properties} emplearon composites centrados en el ciclón para caracterizar la estructura de los ETCs durante su génesis, mostrando que los mecanismos dominantes de desarrollo difieren según la región de formación.

\subsection{PDF de magnitudes y de posición}\label{subsec:tb_kde}

Para caracterizar la distribución estadística de las magnitudes máximas sin imponer supuestos distribucionales a priori, se adopta la Estimación de Densidad por Kernel (KDE) unidimensional, que transforma las observaciones discretas de magnitud $\{x_i\}$ en una función de densidad de probabilidad continua $\hat{f}(x)$, permitiendo contrastar cuantitativamente las distribuciones de precipitación y viento.
Este enfoque evita el sesgo inherente al agrupamiento en clases de un histograma, preservando la ubicación exacta de cada observación, como señaló \citet{weglarczyk2018kernel}.

La selección del ancho de banda $h$ constituye el parámetro crítico del estimador. Se adopta la regla
de Scott, una técnica de selección de ancho de banda ampliamente conocida y también empleada por \citet{weglarczyk2018kernel} entre los métodos de tipo \textit{rule-of-thumb} disponibles para este propósito,
\begin{equation}\label{eq:scott_rule}
h = \hat{\sigma} \cdot m^{-1/5},
\end{equation}
donde $\hat{\sigma}$ es la desviación estándar muestral de las magnitudes.

Formalmente, el estimador unidimensional se define como:
\begin{equation}\label{eq:tb_kde_pdf}
\hat{f}(x) = \frac{1}{m h} \sum_{i=1}^{m} K\!\left(\frac{x - x_i}{h}\right),
\end{equation}
donde $K(\cdot)$ es un núcleo gaussiano estándar:
\begin{equation}\label{eq:tb_gaussian_kernel}
K(u) = \frac{1}{\sqrt{2\pi}} \exp\!\left(-\frac{u^2}{2}\right).
\end{equation}

De forma complementaria, la extensión bidimensional del estimador, denominada aquí PDFe, opera sobre las posiciones relativas $\mathbf{s}_i = (\Delta x_i, \Delta y_i)$ y produce un campo continuo de densidad de probabilidad espacial:
\begin{equation}\label{eq:tb_kde_simple}
\hat{f}(\mathbf{s}) = \frac{1}{m} \sum_{i=1}^{m} K_h(\mathbf{s} - \mathbf{s}_i),
\end{equation}
con núcleo gaussiano espacial:
\begin{equation}\label{eq:tb_gauss_kernel_2d_simple}
K_h(\mathbf{u}) = \frac{1}{2\pi h^2} \exp\!\left(-\frac{\|\mathbf{u}\|^2}{2h^2}\right).
\end{equation}

La moda del campo $\hat{f}(\mathbf{s})$ indica la posición relativa al centro del ciclón donde
estadísticamente es más frecuente encontrar el máximo de la variable.
Los niveles de contorno adoptados para la visualización y la implementación computacional de ambos estimadores se detallan en la Sección~\ref{subsec:prototipos}.

\subsection{Descomposición modal EOF y MEOF}\label{subsec:tb_eof}
El análisis de Funciones Ortogonales Empíricas (EOF) descompone la variabilidad de un campo en modos ortogonales de varianza máxima, separando la señal física del ruido o variabilidad de menor escala \citep{hannachi2007empirical}.
Esta técnica admite extensiones para tratar simultáneamente múltiples grupos de datos: por un lado, el EOF común, orientado a comparar la estructura compartida de un mismo campo entre distintos modelos o reanálisis \citep{hannachi2023eof}; por otro, el EOF combinado o multivariado (MEOF), que concatena campos físicamente distintos en un único vector de estado para capturar su covariabilidad conjunta \citep{liang2018multivariate}. Para aislar los patrones espaciales coherentes del acoplamiento físico entre viento y precipitación, este estudio adopta esta segunda variante, que llamaremos MEOF.

Dado un conjunto de $m$ campos anomalía $X'(\tau, \mathbf{s})$ distribuidos en una grilla espacial de $p$ puntos, se organizan en una matriz $\mathbf{X}$ de dimensiones $m \times p$. La matriz de covarianza espacial $\mathbf{S} = \frac{1}{m-1}\mathbf{X}^T\mathbf{X}$ se descompone espectralmente:
\begin{equation}\label{eq:eof_decomp}
\mathbf{S}\,\mathbf{e}_k = \lambda_k \mathbf{e}_k,
\end{equation}
donde $\mathbf{e}_k$ es el $k$-ésimo autovector (el patrón espacial $\text{EOF}_k$) y $\lambda_k$ su autovalor asociado (fracción de varianza explicada).
Las amplitudes de proyección de cada evento sobre el modo dominante, las componentes principales PC$_k(\tau) = \mathbf{X}(\tau)\,\mathbf{e}_k$, constituyen un indicador del grado de ajuste de cada evento al patrón estructural arquetípico.
Una distribución de PC$_1$ con asimetría positiva o leptocurtosis sugiere la existencia de un subconjunto de sistemas que se proyecta con especial intensidad sobre el modo dominante; una correlación significativa entre PC$_1$ y la vorticidad relativa a 850 hPa asociada a la fase indicaría que el lado del patrón dominante en que cae cada ciclón está sistemáticamente asociado a su intensidad (el signo de esa correlación depende de la convención arbitraria del EOF calculado).
Estas propiedades se evalúan empíricamente en la Sección~\ref{subsec:analisis_pc}.

En el caso multivariado (MEOF), los campos normalizados de precipitación y viento se concatenan en un vector de estado combinado $\mathbf{Y}(\tau) = [Z_1(\tau), Z_2(\tau), \dots, Z_N(\tau)]$ de longitud $Np$, siguiendo el enfoque de PCA combinada (CPCA) formalizado por \citet{bretherton1992intercomparison}.
La misma descomposición espectral se aplica a la matriz de covarianza de $\mathbf{Y}$, obteniendo modos acoplados que describen la covariabilidad conjunta de las variables.
\citet{bretherton1992intercomparison} advirtieron que este enfoque puede sesgarse hacia el EOF dominante de un único campo cuando la razón señal--ruido es baja o el tamaño muestral es reducido, favoreciendo un modo que en realidad refleja la variabilidad individual de precipitación o de viento por separado en lugar de la señal de acoplamiento genuina. 
Por ello, la significancia estadística de los patrones MEOF obtenidos se somete al remuestreo no paramétrico descrito en la Sección~\ref{subsec:tb_bootstrap} antes de interpretarlos como estructuras físicamente robustas.
El primer modo multivariado (MEOF$_1$) produce simultáneamente un patrón espacial para precipitación ($\text{EOF}_{1,P}$) y otro para viento ($\text{EOF}_{1,W}$), ambos extraídos del mismo modo dominante de covariabilidad conjunta.
La especificación formal de la descomposición, el procedimiento de normalización y los criterios de umbralización de contornos se detallan en las Secciones~\ref{subsec:eof_univariado} y~\ref{subsec:eof_multivariado}.

\subsection{Bootstrap-t}\label{subsec:tb_bootstrap}

Para evaluar la robustez estadística de los patrones espaciales derivados del análisis EOF y aislar las señales físicas del ruido estocástico, se adopta el método de remuestreo no paramétrico Bootstrap-t.
A diferencia de las pruebas paramétricas estándar, que asumen distribuciones normales, el Bootstrap-t no impone supuestos distribucionales, lo que lo hace apropiado para las distribuciones de precipitación y viento ciclónico, intrínsecamente asimétricas y de cola pesada.

Esta variante corrige dicha asimetría al estandarizar cada muestra de remuestreo por su propio error estándar, lo que estabiliza la distribución del estadístico frente a la forma de la población y produce intervalos de confianza con cobertura más precisa que el Bootstrap percentílico estándar, como estableció \citet{hesterberg2015bootstrap}.
Si $\hat{\theta}$ es el estimador original (carga del EOF en un punto de grilla) y $\hat{\theta}^*_b$ su réplica en la $b$-ésima muestra bootstrap, el estadístico pivotal se define como
\begin{equation}\label{eq:bootstrap_t}
t^*_b = \frac{\hat{\theta}^*_b - \hat{\theta}}{\widehat{SE}^*_b},
\end{equation}
donde $\widehat{SE}^*_b$ es el error estándar de la submuestra. El intervalo de confianza con nivel $1-\alpha$ se obtiene a partir de los cuantiles empíricos $q_{\alpha/2}$ y $q_{1-\alpha/2}$ de la distribución de $t^*$:
\begin{equation}\label{eq:bootstrap_ci}
\left(\hat{\theta} - q_{1-\alpha/2}\,\widehat{SE},\; \hat{\theta} - q_{\alpha/2}\,\widehat{SE}\right).
\end{equation}
Los puntos de grilla se consideran estadísticamente significativos si su intervalo de confianza excluye el cero. La especificación del número de iteraciones y la implementación se detallan en la Sección~\ref{subsec:bootstrap_significancia}.
